\begin{pastexam}[label={px:gk3-109-ds-03}]{高考三級資訊處理 109 年 第三題}
（申論；25 分）

請回答下列關於 AVL 樹（AVL Tree）的問題：
\begin{itemize}[leftmargin=2em, labelsep=0.3em, itemsep=2pt]
\item[㈠] 我們欲將所管理的鍵值（Key）依序列出，請問是否可以利用一個 AVL 樹對鍵值來進行排序
（Sorting）？若不行，請說明原因；如果可以，請描述方法及時間複雜度。（5 分）
\item[㈡] 請提供一個線性時間的演算法來判斷一個二元搜尋樹是否為 AVL 樹。（10 分）
\item[㈢] 在 AVL 樹上進行一個加入（Insert）操作後，是否最多只需要一次的重構（Restructuring）即可
恢復其平衡的特性？請說明原因。（10 分）
\end{itemize}
\end{pastexam}

\begin{pastsoln}
\textbf{㈠ 可以。} AVL 樹是二元搜尋樹，中序走訪會由小到大列出所有鍵（資結書定理~9.1）。 方法：
\begin{enumerate}[nosep]
\item 把 $n$ 個鍵逐一插入一棵空的 AVL 樹。 AVL 樹高度 $O(\lg n)$（資結書定理~10.1），每次插入（含
旋轉）$O(\lg n)$，共 $O(n \lg n)$。
\item 對樹做中序走訪，依序輸出，每個節點拜訪一次，$O(n)$。
\end{enumerate}
總時間 $O(n \lg n)$，和比較排序的下界 $\Omega(n \lg n)$ 同級，是最佳的；額外空間 $O(n)$（$n$ 個
節點）。 若鍵本來就存在 AVL 樹裡，只要第 2 步，$O(n)$ 就能依序列出。 鍵可能重複時，約定相等的鍵往
右子樹放（或在節點上記次數）即可。

\textbf{㈡ 線性時間判斷。} 題目已說是二元搜尋樹，只要檢查「每個節點的左右子樹高度差 $\le 1$」。
逐個節點分別算高度會重複計算（最壞 $\Theta(n^2)$）；改用\textbf{後序走訪}，先算出左右子樹的高度，
再算自己的，每個節點只處理一次。 用 $-2$ 當「這棵子樹已經不平衡」的記號（空樹高度 $-1$）。
\begin{pseudo}{判斷二元搜尋樹是否為 AVL 樹}
\Function{Check}{$x$} \Comment{平衡時傳回高度，不平衡時傳回 $-2$}
    \If{$x = \NIL$}
        \State \Return $-1$
    \EndIf
    \State $h_L \gets$ \Call{Check}{$x.\mathit{left}$}
    \If{$h_L = -2$}
        \State \Return $-2$
    \EndIf
    \State $h_R \gets$ \Call{Check}{$x.\mathit{right}$}
    \If{$h_R = -2$ \OR $|h_L - h_R| > 1$}
        \State \Return $-2$
    \EndIf
    \State \Return $1 + \max(h_L, h_R)$
\EndFunction
\Function{Is-AVL}{$T$}
    \State \Return (\Call{Check}{$T.\mathit{root}$} $\ne -2$)
\EndFunction
\end{pseudo}
\textbf{正確性}（對子樹大小歸納）：\proc{Check}($x$) 傳回 $-2$ 當且僅當以 $x$ 為根的子樹裡有某個節點
左右高度差 $> 1$；否則傳回它的高度。 空樹成立；非空時，若左或右子樹已不平衡就傳回 $-2$，否則
兩邊的傳回值就是高度，再檢查 $x$ 自己的高度差，正確。 \textbf{時間}：每個節點恰好被呼叫一次
（加上 $n + 1$ 次 \NIL），每次常數時間，$\Theta(n)$。 若也不確定它是不是二元搜尋樹，可以另做一次
中序走訪，檢查輸出是否遞增，同樣 $O(n)$。

\textbf{㈢ 是，最多一次重構}（一次單旋轉或一次雙旋轉）。

理由（資結書第~10.2.3~節、定理~10.2）：插入新節點只會讓「新節點到根的路徑」上的節點高度改變，
所以只有這些祖先可能失衡。 由下往上找到\textbf{第一個}失衡的節點 $z$，它的平衡因子變成 $\pm 2$。
\begin{itemize}[nosep]
\item 插入前 $z$ 子樹的高度設為 $h + 2$，新節點落在較高的那一側（否則 $z$ 不會失衡），使那一側
從 $h + 1$ 長成 $h + 2$，$z$ 子樹變成 $h + 3$。
\item 依新節點在 $z$ 的哪個孫子底下，分 LL、RR（單旋轉）和 LR、RL（雙旋轉）四種。 逐一檢查各
子樹的高度可知，旋轉之後新的子樹根的兩邊高度都是 $h + 1$，子樹高度回到 \textbf{$h + 2$，和插入前
一樣}。
\item 這棵子樹以外的節點，看到的子樹高度和插入前完全相同，插入前它們是平衡的，現在也一樣平衡，
不需要再處理。
\end{itemize}
因此一次插入最多做一次重構。 例：依序插入 $3, 2, 1$，插入 $1$ 後 $3$ 的平衡因子是 $+2$（LL 型），
對 $3$ 右旋得到根 $2$、左 $1$、右 $3$。 這棵子樹插入前（$3$ 帶左子 $2$）是兩層，旋轉後仍是兩層，
上面若還有祖先也不受影響。

對照：\textbf{刪除}就不一樣了，旋轉後子樹高度可能比刪除前少一，失衡會往上傳，最壞要在路徑上
每一層各重構一次，$O(\lg n)$ 次（資結書第~10.2.6~節）。
\end{pastsoln}
